\section{Processing chain}
\label{sec:processing}

\subsection{Overview}
\label{sec:proc-overview}

Raw 20 hertz records pass through eleven stages before a flux is written to the
archive. Figure~\ref{fig:pipeline} shows the chain. Every stage is versioned and
every record carries the version identifier of the chain that produced it, so a
reprocessing campaign can be distinguished from an instrument change in the
archive without reference to an external log.

\begin{figure}[t]
\centering
\begin{tikzpicture}[
  >={Stealth[length=2.2mm]},
  box/.style={draw,rounded corners=2pt,minimum width=24mm,minimum height=8.5mm,
              align=center,font=\scriptsize},
  store/.style={draw,cylinder,shape border rotate=90,aspect=0.22,
                minimum width=17mm,minimum height=11mm,align=center,font=\scriptsize},
  lab/.style={font=\scriptsize,midway,above,inner sep=1pt}
]
  \node[store] (raw) {raw\\ 20 Hz};
  \node[box,right=9mm of raw] (despike) {despike\\ and range};
  \node[box,right=7mm of despike] (rotate) {coordinate\\ rotation};
  \node[box,right=7mm of rotate] (lag) {lag\\ maximisation};
  \node[box,below=9mm of lag] (block) {block average\\ and detrend};
  \node[box,left=7mm of block] (cov) {covariance\\ 30 min};
  \node[box,left=7mm of cov] (wpl) {density\\ correction};
  \node[box,below=9mm of wpl] (spec) {spectral\\ correction};
  \node[box,right=7mm of spec] (foot) {footprint\\ model};
  \node[box,right=7mm of foot] (qc) {quality\\ flags};
  \node[store,right=9mm of qc] (arch) {archive\\ 30 min};

  \draw[->,thick] (raw) -- (despike);
  \draw[->,thick] (despike) -- (rotate);
  \draw[->,thick] (rotate) -- (lag);
  \draw[->,thick] (lag) -- (block);
  \draw[->,thick] (block) -- (cov);
  \draw[->,thick] (cov) -- (wpl);
  \draw[->,thick] (wpl) -- (spec);
  \draw[->,thick] (spec) -- (foot);
  \draw[->,thick] (foot) -- (qc);
  \draw[->,thick] (qc) -- (arch);

  \node[box,above=9mm of foot,fill=white] (slow) {slow suite\\ 1 min};
  \draw[->,thick,dashed] (slow) -- (wpl);
  \draw[->,thick,dashed] (slow) -- (foot);

  \begin{scope}[on background layer]
    \node[draw,dashed,inner sep=2.5mm,fit=(despike)(rotate)(lag)(block)(cov)] (stage1) {};
  \end{scope}
  \node[font=\scriptsize,above=0.5mm of stage1.north west,anchor=south west]
    {high frequency stages, Section~\ref{sec:proc-highfreq}};
\end{tikzpicture}
\caption{Processing chain from raw high frequency records to the archived half
hourly flux. Dashed inputs come from the slow response suite. The footprint
model consumes both the rotated wind statistics and the stability parameter
derived from the slow suite.}
\label{fig:pipeline}
\end{figure}

\subsection{High frequency stages}
\label{sec:proc-highfreq}

Despiking removes points more than 5.5 median absolute deviations from a 5
minute running median, with a maximum of four consecutive points removed before
the whole averaging interval is rejected. Coordinate rotation uses the planar
fit method applied per site and per wind sector of 30 degrees, with the fit
coefficients recomputed monthly.

The turbulent flux of a scalar $c$ over an averaging interval is the covariance
\begin{equation}
  F_c = \overline{w' \rho_c'},
  \label{eq:covariance}
\end{equation}
where $w$ is the vertical velocity, $\rho_c$ the scalar density, primes denote
deviations from the interval mean, and the overbar is the interval average. The
sensible heat flux follows as $H = \rho c_p \overline{w'T'}$ and the latent heat
flux as $\lambda E = \lambda \overline{w' \rho_v'}$.

Because the open path analyser measures density rather than mixing ratio, the
flux requires the correction for expansion and contraction of the air parcel
\citep{duvall2015density},
\begin{equation}
  F_c^{\mathrm{corr}} = F_c
    + \mu \frac{\overline{\rho_c}}{\overline{\rho_a}} \overline{w'\rho_v'}
    + \left( 1 + \mu \frac{\overline{\rho_v}}{\overline{\rho_a}} \right)
      \frac{\overline{\rho_c}}{\overline{T}} \overline{w'T'},
  \label{eq:wpl}
\end{equation}
with $\mu$ the ratio of the molar masses of dry air and water vapour. The
correction is a substantial fraction of the carbon flux at this network: its
median magnitude over the year is 19 percent of the uncorrected flux, and at the
saltmarsh sites in summer it exceeds 40 percent.

\subsection{Spectral corrections}
\label{sec:proc-spectral}

High frequency attenuation from sensor separation, path averaging, and the
finite response of the analyser is corrected by an empirical transfer function
fitted to the site cospectra \citep{grothe2018spectral}. The correction factor is
\begin{equation}
  \Phi = \frac{\int_0^{\infty} C_{wc}^{\mathrm{ref}}(f) \, df}
              {\int_0^{\infty} C_{wc}^{\mathrm{ref}}(f) \, T(f) \, df},
  \label{eq:spectral-correction}
\end{equation}
where $C_{wc}^{\mathrm{ref}}$ is a reference cospectral shape and $T(f)$ the
fitted transfer function. Low frequency loss from the 30 minute averaging window
is assessed by the ogive test and found to be below 3 percent at all sites
except MC08, whose 8 metre height and elevated position give it the largest
transporting eddies in the network.

\subsection{Footprint model}
\label{sec:proc-footprint}

Source area is estimated with a parameterised two dimensional model
\citep{lindholm2017footprint} driven by
the friction velocity, the Obukhov length, the standard deviation of the lateral
wind, and the roughness length. The output per record is the distance containing
80 percent of the flux source, and the fraction of the source area falling
within the intended cover class.

Records with less than 70 percent of the source area in the intended class are
flagged. Across the network this flag affects 6.2 percent of records, rising to
19.8 percent at MC06 for the reason given in Section~\ref{sec:network-siting}.

\subsection{Gap filling}
\label{sec:proc-gapfill}

Gaps are filled by marginal distribution sampling \citep{kaale2022gapfill} on a window of plus or minus
seven days, conditioned on incoming shortwave radiation, air temperature, and
vapour pressure deficit, with a fallback to the mean diurnal course when no
comparable meteorological conditions exist in the window. Gaps longer than 12
days are not filled and the affected period is excluded from the annual sums with
a corresponding widening of the reported uncertainty.

The uncertainty contributed by gap filling is estimated by artificially removing
observed records at the same rate and in the same diurnal and seasonal pattern
as the real gaps, then comparing the filled value with the withheld observation.
The resulting contribution to the annual budget is 17 grams of carbon per square
metre at the saltmarsh sites and 11 at the dune sites.
